\subsubsection[多变量函数的可微性]{Differentiability}

Recall: $f$ is differentiable at $x_0$ iff $f(x_0+\Delta x)-f(x_0)=a\Delta x+\circ\l(\lvert\Delta x\rvert\r),\ \Delta x\to 0$

($\iff f$ has derivative at $x_0$ and $a=f'(x_0)$, $\lim\limits_{\Delta x\to 0}\dfrac{f(x_0+\Delta x)-f(x_0)}{\Delta x}$ exist)

\begin{definition}
    Let $f:\,D(\subseteq\mathbb{R}^n)\to\mathbb{R}$, $x$ is an interior point of $D$ if
    $$
        f(x+h)=f(x)+A(h)+\circ\l(\lvert h\rvert\r),\ h\to0
    $$
    \emph{$\implies$ If $f$ is differentiable at $x$, then the graph of $f$ has a tangent \emph{hyperplane} (in $\mathbb{R}^{n+1}$) which is the graph of $\substack{h\mapsto f(x)+A(h)\\\mathbb{R}^n\to\mathbb{R}}$}\\
    When $A:\,\mathbb{R}^n\to\mathbb{R}$ which is a linear function, then we say $f$ is \emph{defferentiable} at $x$, the map $A$ is called the differential of $f$ at $x$, denoted by $df|_x$

    If $f$ is differentiable at any $x\in D$, then we say $f$ is differentiable on $D$ and $\{df|_x\mid x\in D\}=df$ is called the differential of $f$ on $D$.
\end{definition}

\begin{remark}
    For $n=1$, $\{f'(x)\mid x\in D\}=f'$ is called the derivative function of $f$ on $D$. We also have $df=f'(x)dx$, called the differential of $f$ on $D$.
\end{remark}

\begin{theorem}
    If $f$ is differentiable at $x$, then it is continuous at $x$.
\end{theorem}

Assume $f$ is differentiable at $x$, then $f(x+h)=f(x)+A(h)+\circ\l(\lvert h\rvert\r),\ h\to0$

$\implies f(x+te_i)=f(x)+A(te_i)+\circ\l(\lvert t\rvert\r)=f(x)+tA(e_i)+\circ\l(\lvert t\rvert\r),\ t\to0$ ($i=1,2,\dots,n$)

$\implies \dfrac{\partial f}{\partial x_i}(x)$ exists (since $\dfrac{\partial f}{\partial x_i}(x)=\lim\limits_{t\to 0}\dfrac{f(x+te_i)-f(x)}{t}$), and $\dfrac{\partial f}{\partial x_i}(x)=A(e_i)$

\begin{theorem}
    If $f$ is differentiable at $x$, then $\dfrac{\partial f}{\partial x_1}(x),\dots,\dfrac{\partial f}{\partial x_n}(x)$ all exist, and
    $$
        \dfrac{\partial f}{\partial x_i}(x)=df|_x(e_i)
    $$
    \textbf{The converse is not true.}
\end{theorem}

Look at some special functions: $\boxed{X_i:\,\mathbb{R}^n\to\mathbb{R},\,i=1,\dots,n}$ (The $n$ Cartesian coordinates function)

\allbold{$x_i(x+h)=x_i(x)+h_i$, where $h_i=i$-th component of $h$ under the standard basis $\{e_1,\dots,e_n\}$}

\allbold{$\implies x_i$ is differentiable on $\mathbb{R}^n$, so we have $dx_i$}

\begin{align*}
    f(x+h) & =f(x)+df|_x(h)+\circ(|h|)                                                                \\
           & =f(x)+df|_x\!\left(\sum_{i=1}^{n}h_ie_i\right)+\circ\l(\lvert h\rvert\r)                 \\
           & =f(x)+\sum_{i=1}^{n}df|_x(e_i)h_i+\circ\l(\lvert h\rvert\r)                              \\
           & =f(x)+\sum_{i=1}^{n}\dfrac{\partial f}{\partial x_i}(x)h_i+\circ\l(\lvert h\rvert\r)     \\
           & =f(x)+\sum_{i=1}^{n}\dfrac{\partial f}{\partial x_i}(x)dx^i(h)+\circ\l(\lvert h\rvert\r)
\end{align*}

$$
    f(x)+\sum_{i=1}^{n}\dfrac{\partial f}{\partial x_i}(x)h_i+\circ\l(\lvert h\rvert\r)
$$
\allbold{$\rightarrow df|_x=\sum_{i=1}^n\dfrac{\partial f}{\partial x_i}(x)\,dx_i$} ``The linear \emph{combination} of the $n$ special linear function $dx_1,\dots,dx_n$, with coefficients the partial derivatives.''

\emph{$dx_1,\dots,dx_n$ are the standard differential 1-form.}

\allbold{Further, we can write}

$$df=\sum_{i=1}^{n}\dfrac{\partial f}{\partial x_i}\,dx_i$$

``module'' 模

$$\left(\dfrac{\partial f}{\partial x_1},\dots,\dfrac{\partial f}{\partial x_n}\right)$$

\allbold{$\operatorname{Grad}f|_x=\nabla f|_x=\sum_{i=1}^{n}\dfrac{\partial f}{\partial x_i}(x)e_i$ is called the gradient 梯度 of $f$ at $x$.}
